\uuid{vKlv}
\titre{Limite en l'infini}
\auteur{bodin ; hanani ; mzari}
\organisation{exo7}
\licence{CC BY-NC-SA 4.0}
\datecreate{2020-01-20}
\niveau{L1}
\module{Analyse}
\chapitre{Continuité, limite et étude de fonctions réelles}
\sousChapitre{Limite de fonctions}
\difficulte{3}
\theme{Limites des fonctions réelles, limite en l'infini}
\exo7qcmid{lille-1.limites.030}

\contenu{
\begin{question}
\qclassification{123.03}
\qtags{motcle=limite en l'infini}
%Soit  $f(x)=\frac{(x^x)^x}{x^{(x^x)}}$.
Soit  $f(x)=\frac{(2x)^x}{x^{(2x)}}$. Quelles sont les assertions vraies ?
\begin{answers}

    \bad{$\lim_{x\to +\infty} f(x)=+\infty$}

    \good{$\lim_{x\to +\infty } f(x)=0$}

    \bad{$f$ n'admet pas de limite en $+\infty$.}

    \bad{$\lim_{x\to +\infty} f(x)=1$}

\end{answers}
\begin{explanations}
Par définition, si $u$ et $v$ sont deux fonctions telles  que $u>0$, $u^v=e^{v\ln u}$. On en déduit que $f(x)=\exp[ x\ln (2x )- 2x\ln x] =  \exp[ x\ln 2 - x\ln x]$. Donc $\lim_{x\to +\infty } f(x)=0$.
\end{explanations}
\end{question}
}
